\documentclass[10pt,a4paper]{amsart}
\usepackage[margin=19mm]{geometry}
\usepackage{amsmath,amssymb,mathtools}
\usepackage[scr=rsfs,cal=euler]{mathalfa}
\usepackage{xfrac}
\usepackage[unicode,hidelinks,bookmarks=false]{hyperref}
\newcommand{\ie}{i.e.\ }
\newcommand{\cf}{cf.\ }
\newcommand{\resp}{respectively\ }
\newcommand{\Subsection}{\subsection}
\providecommand{\nobreakdash}{\nobreak\hbox{-}}
\setlength{\emergencystretch}{2em}
\multlinegap0pt
\theoremstyle{plain}
\newtheorem{theorem}{Theorem}
\newtheorem{corollary}[theorem]{Corollary}
\newtheorem{lemma}[theorem]{Lemma}
\theoremstyle{remark}
\newtheorem{remark}{Remark}
\title[Non-Hermitian band matrices]{Outliers and bounded rank perturbation for non-Hermitian random band matrices: a matrix-unit block computation for the model parameters}
\author{Yi HAN}
\date{Source: arXiv:2408.00567v4 (CC BY 4.0)}
\begin{document}
\maketitle
\noindent\emph{Source and licence.} Outliers and bounded rank perturbation for non-Hermitian random band matrices. Version arXiv:2408.00567v4 (CC BY 4.0), distributed under the Creative Commons Attribution 4.0 International licence, \url{http://creativecommons.org/licenses/by/4.0/}, as recorded in the arXiv licence metadata for this exact version and shown on the abstract page \url{https://arxiv.org/abs/2408.00567v4}. Full licence notice: licenses/arxiv-2408.00567-CC-BY-4.0.txt.\par\smallskip

\noindent\textit{Author byline.} The author list is the article's own single-author byline; the arXiv metadata that accompanies this version spells the family name "Han" while the article itself sets it "HAN", and this package uses the article's own form "Yi HAN" so that the delivered byline agrees character for character with the article.\par\smallskip

\noindent\textit{Editorial scope.} Counted unit: the article's Lemma with source label \texttt{lem:block-parameter} (source0.tex lines 1089--1130), the matrix-unit block computation that verifies, for a random matrix presented as a sum of independent zero-mean blocks of matrix units, that the four model parameters of the article and the block supremum bound take the stated values under the article's hypotheses, delivered together with the article's complete original proof of it (source0.tex lines 1133--1192, one \texttt{proof} environment of 60 lines, adjacent to the statement). No mathematics is written, repaired or re-derived: statement and proof are the article's own text line for line, in the article's own notation (the block decomposition with matrix units, the distinctness of the units, the block covariance and pseudo-covariance bounds, the doubly stochastic normalisation, the parameters and the block supremum norm), and no retained line is substituted, re-flowed or omitted. No further part of the article is required as context and none is retained: the counted proof uses only the definitions introduced in its own statement, so no cross-reference and no citation occurs in any retained line; consequently the wrapper carries no bibliography and no bibliography entry, and no context passage is delivered. Nothing else of the article is retained: its main theorems, its other lemmas, its proofs and its appendices are omitted. Editorial formatting only: an AMS \texttt{amsart} preamble that restates the article's theorem-like environment declarations, in which Theorem, Corollary and Lemma share one counter as in the article, and that adds no mathematical macro, because every control sequence used by the retained lines is standard. The article's own source declares the class \texttt{amsart} as well; no class or style file is shipped with this package. Numbering differs from the article, because the article declares its lemma as sharing the theorem counter and resets that counter in each section, while the wrapper keeps the shared counter without a section reset and retains only one environment: the counted lemma therefore prints as Lemma 1, whereas the article prints it as a lemma in its fourth section. No picture environment and no graphical inclusion occurs in any retained line, and the package delivers no image asset. One bundled source file, \texttt{sources/163211/source0.tex}, accompanies the delivered item as the exact source of every retained span. Every retained span is recorded in \texttt{delivery-evidence.json} with method \texttt{exact\_tex}.
\begin{lemma}[A matrix-unit block computation]\label{lem:block-parameter}
Let the matrix $X$ have the following decomposition into independent blocks $Z_\alpha$,
\[
X=\sum_{\alpha\in\mathcal A} Z_\alpha,\qquad
Z_\alpha=\sum_{r=1}^{q_\alpha} \xi_{\alpha r} a_{\alpha r}E_{i_{\alpha r}j_{\alpha r}},
\]
where the matrices \(Z_\alpha\) are mutually independent and have zero mean. Here $a_{\alpha r}\in\mathbb{C}$, $\xi_{\alpha r}$ are complex-valued random variables and $E_{i_{\alpha r}j_{\alpha r}}$ are the matrix units. Assume that the
matrix units \(E_{i_{\alpha r}j_{\alpha r}}\) appearing above are all distinct, that for some $\kappa_N>0$,
\[
\max_{\alpha,r}|a_{\alpha r}|\le \kappa_N,
\]
and that the block covariance and pseudo-covariance matrices satisfy
\[
\left\|\mathbb E\,\xi_\alpha\xi_\alpha^*\right\|\le L,
\qquad
\left\|\mathbb E\,\xi_\alpha\xi_\alpha^t\right\|\le L
\]
for every \(\alpha\), where \(L=O(1)\) and where \(\xi_\alpha=(\xi_{\alpha 1},\ldots,\xi_{\alpha q_\alpha})^t\). If moreover we are in the doubly stochastic case
\[
\mathbb E[XX^*]=\mathbb E[X^*X]=I,
\]
then we have the following parameter upper bounds
\begin{equation}
\sigma(X)=1,\qquad
\sigma_*(X)\le L^{1/2}\kappa_N,\qquad
v(X)\le L^{1/2}\kappa_N,\qquad
\widetilde v(X)\le L^{1/4}\kappa_N^{1/2}.
\end{equation}
In the non-Gaussian case, we have
\begin{equation}
R(X)\le \max_{\alpha\in\mathcal A}\|Z_\alpha\|_\infty.
\label{4.2}
\end{equation}
In particular, if each \(Z_\alpha\) is supported on a matching of matrix units (so that their row indices are distinct and their column indices are distinct), then
\begin{equation}
R(X)\le
\kappa_N
\max_{\alpha\in\mathcal A}
\left\|\max_{1\le r\le q_\alpha}|\xi_{\alpha r}|\right\|_\infty .
\label{4.3}
\end{equation}
\end{lemma}

\begin{proof}
For unit vectors \(u,w\), denote by
\[
b_{\alpha r}:=a_{\alpha r}\overline{u}_{i_{\alpha r}}w_{j_{\alpha r}}.
\]
Then we compute the variance via
\[
\mathbb E|\langle u,Xw\rangle|^2
=
\sum_{\alpha\in\mathcal A}
\mathbb E\left|\sum_{r=1}^{q_\alpha}\xi_{\alpha r}b_{\alpha r}\right|^2
\le
L\sum_{\alpha,r}|b_{\alpha r}|^2.
\]
Since the matrix units are distinct,
\[
\sum_{\alpha,r}|b_{\alpha r}|^2
\le
\kappa_N^2
\sum_{\alpha,r}|u_{i_{\alpha r}}|^2|w_{j_{\alpha r}}|^2
\le
\kappa_N^2
\sum_{i,j}|u_i|^2|w_j|^2
=
\kappa_N^2.
\]
Thus
\[
\sigma_*(X)^2\le L\kappa_N^2.
\]

Similarly, for any matrix \(M\) with \(\|M\|_{\mathrm{HS}}\le1\),
\[
\operatorname{Tr}(XM)
=
\sum_{\alpha,r}\xi_{\alpha r}a_{\alpha r}(M)_{j_{\alpha r}i_{\alpha r}},
\]
and the same computation gives
\[
\mathbb E|\operatorname{Tr}(XM)|^2
\le
L\kappa_N^2\sum_{i,j}|M_{ji}|^2
\le
L\kappa_N^2.
\]
Equivalently, \(v(X)^2\le L\kappa_N^2\). Since
\[
\mathbb E[XX^*]=\mathbb E[X^*X]=I,
\]
we have \(\sigma(X)=1\). Therefore
\[
\widetilde v(X)^2=v(X)\sigma(X)\le L^{1/2}\kappa_N,
\]
which gives
\[
\widetilde v(X)\le L^{1/4}\kappa_N^{1/2}.
\]
Finally, the estimate for \(R(X)\) exactly follows from its definition, being the largest one among the essential
supremum norm of each independent block $Z_\alpha$. 
\end{proof}

\end{document}
